\section{Twenty-four unrestricted integer unknowns}
\label{sec:integer}

The preceding construction can be adapted to a substantially smaller
integer domain. The current collection has equations in 26 unrestricted
integers; the construction below has \textbf{24}, degree \textbf{1936},
and a division-free evaluator with \textbf{167} binary operations. Its
three-equation version uses 162 operations and has maximum degree 1920.
The proof combines the saturated root norm of package~19 with the odd-scale
signed-domain theorem of package~23 \cite{repo19,repo23,oddscale}.
The new work is the separation argument when the half-binomial scale, the
main Pell parameter, and the separating coefficients can be negative.

\subsection{Domain maps and the equation}

For each of the six ports
\[
 J,F,\alpha,z,\eta,\zeta
\]
use a disjoint copy of the polynomial
\[
 \operatorname{Pos}(x,y,u)=1+x^2+y^2+u(u+1).
\]
Its values are positive, and it represents every positive integer:
apply the three-square theorem to $4n-3$, whose representation has one
odd square and two even squares. Supply six further unrestricted integers
$h,T,t,r_i,r_s,r_w$, and put
\begin{equation}\label{eq:integer-domain}
 i=2r_i+1,\qquad s=2r_s+1,\qquad w=r_w^2.
\end{equation}
Thus the number of supplied integers is $6\cdot3+6=24$.
The executable circuit names these last three inputs \texttt{ii},
\texttt{ss}, and \texttt{ww}, and names the merged coordinate $t$
\texttt{r}.

Retain the common abbreviations $q,C,X,Y,E,k,c,a,H,\Delta,R,N_k,N_t$
of the positive construction. In particular,
\begin{align*}
 X&=wq,&Y&=sq^3,&E&=XY,&k&=\eta+\zeta,\\
 c&=kY+\eta,&a&=E+Y,&H&=4a+3,&\Delta&=a^2+H.
\end{align*}
To make the norm definition explicit, set
\begin{align*}
 L&=kXY^2,& A&=1+L(L+k),&G&=1+\Delta c^2,\\
 P&=X+ac,&I&=ic^2,&D&=1+\Delta I^2,\\
 Q&=\Delta^2I^2,&d&=Q-1,&v&=cT,\\
 b&=c+RD,&e&=Q(Dv^2+b^2)-1,&B_1&=2Qvb.
\end{align*}
Here $L$ is an arithmetic abbreviation, unrelated to a compiler cell
length. All these symbols are computed expressions, not new unknowns.

Use the following signed separating weights:
\begin{align}
 M&=1+8(TI)^2+16Ra i^2c^3
     =1+8(ic)I(T^2c+2Ra),\label{eq:signed-weight}\\
 N&=1+8HM,\qquad Z=Ht+NP.\label{eq:signed-weight-N}
\end{align}
The second expression for $M$ is the factored form used by the circuit.
For formal roots $u^2=A$ and $g^2=G$, define
\begin{align}
 F(u,g)&=\bigl[d(Z-HMu-Ng)^2-H^2e\bigr]^2-H^4B_1^2D,
 \label{eq:signed-factor}\\
 \Omega&=\prod_{u=\pm\sqrt A}\prod_{g=\pm\sqrt G}F(u,g).
 \label{eq:signed-norm}
\end{align}
The product is an ordinary polynomial with integer coefficients; its
division-free circuit does not use square roots or a resultant operation.
The requested single equation is
\begin{equation}\label{eq:integer-final}
 \boxed{\;N_kN_t(1+\Omega)-1=0.\;}
\end{equation}
The corresponding system is $N_k=1$, $N_t=1$, $\Omega=0$.
All four fixed coefficient ports must come from one complete authentic
tuple. The new tuple in Section~\ref{sec:coefficients} has the required
interface; package~19's compact Wang943 tuple is a separate valid choice.

\begin{theorem}\label{thm:integer24}
For the coherent fixed coefficients specified in this report, equation
\eqref{eq:integer-final} has an integer zero if and only if the blank
run of K215 halts. The same statement holds for its three-equation system.
Consequently ZFC proves that the absence of integer zeros is equivalent
to $\operatorname{Con}(\mathrm{ZFC})$.
\end{theorem}

\subsection{Bootstrap with signed scales}

The optimized norm evaluator has the form
$\Omega=U^2-\Theta V^2$, where
$\Theta=4Z^2N^2G$. Thus $\Omega$ is a square modulo four.
At a zero of \eqref{eq:integer-final}, all three integer factors are
units. The alternative $1+\Omega=-1$ would require a square to equal
two modulo four. Therefore
\[
 \Omega=0,\qquad N_k=N_t=\epsilon,\qquad\epsilon\in\{1,-1\}.
\]
For the system these conclusions hold directly with $\epsilon=1$.

The weak transport estimate needs only the six guarded positive ports:
\[
 q-F-z(q-1)=1-F-(z-1)(q-1)\leq0.
\]
If $C\leq-1$, then $N_t\leq-(K_0+w)\leq-2$, a contradiction.
Hence $C\geq0$ and $q-F=C+\alpha>0$. The authentic mask bounds give
\begin{equation}\label{eq:integer-mask-bound}
 q\geq16,\qquad q(q^2-1)<R<q^4-2.
\end{equation}
These estimates are valid before reconstructing any root and do not use
$s>0$ \cite{normalized,repo19}.

Since $s$ is odd, $|Y|\geq q^3$. Also $X\geq0$, and $a,c$ have
the same sign as $Y$. Their absolute values satisfy
$|a|\geq q^3$ and $|c|>k(|Y|-1)>1$. Consequently $|H|>1$ and
\[
 \Delta=(a+2)^2-1\geq(q^3-2)^2-1>q^4>R.
\]
The multiplier $i$ is nonzero, so
\begin{equation}\label{eq:integer-pre-root}
 D>\Delta+|c|,\qquad 1<b<Q,\qquad Q>D>1.
\end{equation}
For the upper bound on $b$, use
$Q-b=(\Delta-R)D-\Delta-c>0$. In particular $b$ is positive even
when $c$ is negative. Finally, $Ra c^3>0$, so $M>1$, and
\begin{equation}\label{eq:integer-N-bound}
 \left|\frac NH\right|>7M,\qquad \gcd(N,H)=1.
\end{equation}

\subsection{Recovering roots without assuming their signs}

Because $A\geq1$ and $G>1$, a zero of \eqref{eq:signed-norm}
selects real roots $u,g$ with $F(u,g)=0$. Work in the totally real
Galois field $\mathbb Q(u,g)$, allowing degrees one, two, or four, and put
\[
 y=\frac{Z-HMu-Ng}{H}.
\]
When $B_1\ne0$, define $f=(e-dy^2)/B_1$. Then $f^2=D$ and every
actual conjugate satisfies
$d(y')^2=Q(vf'-b)^2-1$. When $B_1=0$, necessarily $v=0$, and the
same factor directly gives $d(y')^2=Qb^2-1$. In both cases,
\begin{equation}\label{eq:integer-y-bound}
 |y'|<2\bigl(|v|\sqrt D+b\bigr),\qquad
 Z=HMu+Ng+Hy.
\end{equation}
No division by $T$ or $B_1$ is made in the exceptional branch.

First exclude $w=0$. The inequalities $|a|\geq4096$ and $|c|>1$ give
\[
 \sqrt G>\frac{|a||c|}{2},\qquad
 \sqrt D<2|a||i|c^2,\qquad D<4a^2i^2c^4.
\]
Term by term in \eqref{eq:signed-weight}, these imply
\begin{align}
 M\sqrt G
 &>2|c|+4|a||iT||c|^3+8Ra^2i^2c^4\notag\\
 &>2\bigl(|v|\sqrt D+b\bigr).
 \label{eq:integer-boundary-separation}
\end{align}
Here $i^2T^2c^2\geq|iT|$, including $T=0$, and
$b\leq|c|+RD$. If $w=0$, then $A=1$ and $u$ is already integral.
An automorphism flipping $g$ would contradict
\eqref{eq:integer-N-bound}, \eqref{eq:integer-y-bound}, and
\eqref{eq:integer-boundary-separation}. Thus $g$ is rational and
integral, and $y$ is rational.

We use the rational auxiliary Pell saturation lemma: for integers
$S_0>1,V$, a rational solution of
\begin{equation}\label{eq:integer-saturation}
 S_0^2V^2-(S_0^2-1)y^2=1
\end{equation}
is integral. Its proof by ordinary Pell classification and divisibility
of first coordinates does not assume $S_0^2-1$ squarefree
\cite{repo19}. When $B_1\ne0$, the rational $f$ has integral square,
so is integral, and the lemma applies with
$S_0=|\Delta I|$, $V=vf-b$. When $B_1=0$, apply it with $V=-b$.
Thus $y$ is integral in either case.

Set $\gamma=(g-P)/H$. The linear equation gives
$t=Mu+N\gamma+y$. Hence $H\gamma$ and $N\gamma$ are integers;
coprimality in \eqref{eq:integer-N-bound} makes $\gamma$ integral.
At $X=0$ the main norm would now imply
\[
 1=(ac+H\gamma)^2-\Delta c^2
   =H(2ac\gamma+H\gamma^2-c^2),
\]
contradicting $|H|>1$. Therefore $w>0$ and $X\geq q\geq16$.

For $w>0$, we have $\sqrt A>L$, $|a+2|<2|E|$, and
$|c|<2k|Y|$. Dividing \eqref{eq:integer-y-bound} by $\sqrt A$
and using $L=|E|k|Y|$ yields
\begin{align}
 \frac{|y'|}{\sqrt A}
 &<8|TI|+\frac4{|E|}+16R|E|i^2|c|^3\notag\\
 &<1+8(TI)^2+16R|a|i^2|c|^3=M.
 \label{eq:integer-ordinary-separation}
\end{align}
The last step uses $|E|\geq q^4$, $|a|>|E|$, and the integer
inequality $(TI)^2\geq|TI|$.

For either sign of $Y$, the elementary estimates
\[
 \Delta>\tfrac34a^2,\quad |c|>\tfrac34k|Y|,
 \quad G>\tfrac{27}{64}L^2,\quad A<\tfrac32L^2
\]
give $\sqrt G>\sqrt A/2$. An automorphism flipping $g$ in
\eqref{eq:integer-y-bound} would therefore imply
$|N/H|\sqrt G<2M\sqrt A$, whereas
\eqref{eq:integer-N-bound} gives
$|N/H|\sqrt G>\tfrac72M\sqrt A$. Thus $g$ is rational. Once $g$
is fixed, a flip of $u$ contradicts the uniform bound
\eqref{eq:integer-ordinary-separation}. Hence $u$ is rational as well.
Both are integers, and the saturation and coprimality argument above
makes $y,\gamma$ integers. This uses only actual field automorphisms;
square radicands or coincident quadratic extensions cause no exception.

If $T=0$, the recovered integral $y$ would satisfy
\[
 (Q-1)(y^2-b^2)=b^2-1>0.
\]
Since $1<b<Q$, one has $|y|\geq b+1$ and $Q-1\geq b$. The left
side is at least $b(2b+1)>b^2-1$, a contradiction. Thus $T\ne0$,
and the recovered $f=(e-dy^2)/B_1$ is an integer.

\subsection{Soundness and completeness of the signed domain}

The roots just reconstructed satisfy exactly
\begin{align*}
 u^2&=1+L(L+k),\\
 (X+ac+H\gamma)^2-\Delta c^2&=1,\\
 f^2-\Delta i^2c^4&=1,\\
 Q[c(Tf-1)-Rf^2]^2-(Q-1)y^2&=1,
 \qquad N_k=N_t=\epsilon.
\end{align*}
There are six positive guarded ports, odd $s,i$, nonzero square $w$,
and unrestricted integer roots. These are precisely the hypotheses of
the signed odd-scale theorem \cite{oddscale}. Its negative-$s$ argument
recovers the main Pell rank $R$ and rules out both main-root signs by
their congruences modulo $H$. Hence $s>0$. The inherited signed compiler
then forces $\epsilon=1$ and recovers an accepting blank K215 tableau
\cite{normalized}. Our root recovery did not assume these conclusions.

Conversely, an accepting computation has canonical normalized witnesses
with odd $i$ and square $w$. The additional oddness of $s$ follows from
the native population identities. If $q=2^j$, the native masks give
\[
 \operatorname{popcount}(R)=3j+2,
 \qquad v_2(Y)=\operatorname{popcount}(R)-2=3j.
\]
Since $Y=sq^3$, the scale $s$ is odd. The required mask hypotheses are
$M_C\equiv2\pmod4$, $M_F^{\rm native}\equiv4\pmod8$, and
$\operatorname{popcount}(M_C)+
\operatorname{popcount}(M_F^{\rm native})=d_{\rm base}=\log_2B$;
the coefficient adapter preserves them. The arithmetic port is still
the \emph{shifted} $M_F$.

Choose $t=y+Mu+N\gamma$. Then one norm factor vanishes and
$N_k=N_t=1$. The six positive coordinates admit their three-integer
representations. This proves Theorem~\ref{thm:integer24} after composition
with the inherited K215/refutation calibration.

\subsection{Exact degree and operation count}

Assign weight two to $J,F,\alpha,z,\eta,\zeta,w$ and weight one to
$s,i,h,T,t$. The principal degrees are
\begin{align*}
 \deg(a,H,c,\Delta,L,A,G)&=(11,11,9,22,20,40,40),\\
 \deg(R,I,D,Q,d,b,e,B_1)&=(8,19,60,82,82,68,218,160),\\
 \deg(M,N,Z,N_k,N_t)&=(48,59,79,12,4).
\end{align*}
Write a subscript $0$ for the highest homogeneous part. Then
$A_0=G_0=L_0^2$, $P_0=L_0=a_0c_0$, and
\[
 M_0L_0=16b_0,\qquad N_0=8H_0M_0.
\]
In the four leading sign choices, $Z-HMu-Ng$ has leading part
$H_0M_0L_0$ times one of $-1,1,15,17$. Thus the leading coefficients
of the inner bracket in \eqref{eq:signed-factor}, after extracting
$H_0^2d_0b_0^2$, are
\[
 255,\quad255,\quad57599,\quad73983.
\]
Their squares have degree 480, while the competing term $H^4B_1^2D$
has degree 424. Therefore
\begin{equation}\label{eq:integer-top-form}
 \Omega_0=C_*H_0^{16}d_0^8b_0^{16},\qquad
 C_*=255^4(57599\cdot73983)^2>0.
\end{equation}
Its degree is $16\cdot11+8\cdot82+16\cdot68=1920$.
Multiplying $N_kN_t$ adds $12+4$, giving exact degree 1936.
The squared-residual alternative has degree 3840.

For an explicit uniform certificate, in the actual circuit variable
names the coefficient of
\[
 J_x^{1354}h\,\mathrm{ii}^{48}\eta_x^{192}
 \mathrm{ww}^{162}\mathrm{ss}^{177}z_x^2
\]
is $2^{257}C_*\beta^{677}$. The exponent sum is 1936 and the
coefficient is nonzero for every authentic $\beta>0$. Thus the degree
claim applies to the full fixed integers, independently of numerical
specializations.

The optimized positive core takes 120 product operations; the signed
weight takes two more. The domain maps initially cost 47 operations.
Two additions disappear by absorbing the constant terms of $F$ and $z$:
if $F=1+F_*$, $z=1+z_*$, and $r=\beta J$, then
\[
 q-F=r-F_*,\qquad N_t=(K_0+w)C-F_*-z_*r.
\]
The final counts are therefore 167 operations for the product, 162 for
the system, and 170 for the sum of squares. All supplied coordinates and
all emitted gates are live. Exact conjugate-product checks, domain-map
identity checks, and independent finite-field leading-coefficient checks
corroborate the circuit. The universal equivalence rests on the proof
above and the explicitly stated inherited dependencies, rather than on
those finite diagnostics.
